\documentclass[12pt,a4paper]{amsart}

\usepackage{amsmath,amssymb,amsthm}
\usepackage{mathtools}
\usepackage[colorlinks=true,linkcolor=blue!60!black,citecolor=green!50!black,
            urlcolor=blue!70!black]{hyperref}
\usepackage[shortlabels]{enumitem}
\usepackage{booktabs}

%% ---- Theorem environments ----------------------------------------
\newtheorem{theorem}{Theorem}[section]
\newtheorem{lemma}[theorem]{Lemma}
\newtheorem{proposition}[theorem]{Proposition}
\newtheorem{corollary}[theorem]{Corollary}
\theoremstyle{definition}
\newtheorem{definition}[theorem]{Definition}
\newtheorem{remark}[theorem]{Remark}
\newtheorem{problem}{Problem}[section]

%% ---- Notation ----------------------------------------------------
\newcommand{\PW}{\mathrm{PW}_\omega}
\newcommand{\Kop}{\mathcal{K}_c}
\newcommand{\Kul}{\kappa_c}
\newcommand{\GVM}{G^{(N)}_{\mathbf{p}}}
\newcommand{\AM}{\alpha_N}
\newcommand{\BM}{\beta_N}
\newcommand{\QM}{Q^{(N)}}
\newcommand{\psin}{\psi_n^{(c)}}
\newcommand{\lam}{\lambda_n(c)}
\newcommand{\Tmin}{T_{\min}}
\newcommand{\EN}{E_N}
\newcommand{\norm}[1]{\|#1\|}
\newcommand{\ip}[2]{\langle #1,\, #2\rangle}

%% ---- Fourier convention ------------------------------------------
%% Engineering: \hat f(\xi)=\int f(x)e^{-2\pi ix\xi}dx.
%% Consistent with Bonami--Karoui. Plancherel: \|f\|_{L^2}=\|\hat f\|_{L^2}.

\title[Coercivity of the Prolate Gram Form]
      {Coercivity of the Prolate Gram Form at Prime Sampling Points}

\thanks{This paper is the first in a series on spectral
constructions related to the Hilbert--P\'{o}lya conjecture.
Paper~II develops the scaling limit and trace formula;
Papers~III--VI address phase obstruction, no-go results,
and Mellin bandwidth decay.
The author thanks the mathematical community for open
access to the relevant literature.}

\author{Anonymous Author}
\address{}
\email{}
\date{\today}

\subjclass[2020]{%
  42C05,  % Orthogonal functions and polynomials
  42B10,  % Fourier and Fourier--Stieltjes transforms
  47B32,  % Operators on Hilbert spaces
  11M06   % Zeta and $L$-functions: analytic theory
}

\keywords{prolate spheroidal wave functions, Gram form, coercivity,
Kulikov constant, prime sampling, Weil positivity, Paley--Wiener space,
Connes--Consani--Moscovici}

\begin{document}
\maketitle

\begin{abstract}
Let $\PW$ denote the Paley--Wiener space of $\omega$-bandlimited
functions, and let $\{\psin\}_{n=0}^{N-1}$ be the prolate spheroidal
wave functions (PSWFs) forming an orthonormal basis of the
$N$-dimensional prolate subspace $V_N\subset\PW$, with
time-bandwidth product $c=\omega T$.
For $N$ distinct sampling points $\mathbf{p}=(p_1,\ldots,p_N)\subset[-T,T]$
we define the \emph{prolate Gram form}
$\QM(f)=\sum_{j=1}^N|f(p_j)|^2$, $f\in V_N$.

We prove that $\QM$ is \emph{coercive}:
there exists $\AM>0$ (depending only on $N$, $c$, $\mathbf{p}$) such that
\[
  \QM(f)\;\geq\;\AM^2\,\norm{f}_{L^2}^2
  \qquad\text{for all }f\in V_N.
\]
The proof is \emph{unconditional}: coercivity follows from the
Chebyshev-system structure of $V_N$ (Lemma~\ref{lem:B}) and requires
no quantitative sampling assumption.
Under the additional Kulikov--Slepian condition
$N\leq(2/\pi-\varepsilon)c$ and a sampling-quality assumption
\textup{(T.2)} on~$\mathbf{p}$, we further obtain an explicit
exponential lower bound $\AM\geq\AM^{\mathrm{exp}}>0$, spectral
stability, a Weil-type decomposition, and a limit statement as
$c\to\infty$.

The result provides a rigorous finite-dimensional building block
for the Weil positivity program of Connes--Consani--Moscovici
\cite{CCM2025}. Applications to prime sampling points are
highlighted; the passage to $N\to\infty$ remains open.
\end{abstract}

\tableofcontents

%%% =================================================================
\section{Introduction}
%%% =================================================================

The \emph{Weil positivity criterion} \cite{Weil1952,Bombieri2000,CC2019}
states that the Riemann Hypothesis is equivalent to positivity of
\[
  W(h)\;=\;
  \sum_{\substack{p\,\text{prime}\\k\geq1}}
  \log p\cdot\bigl|\widehat{h}(\log p^k)\bigr|^2
  +\text{(archimedean terms)}
\]
for all admissible test functions $h$.
The program of Connes--Consani--Moscovici \cite{CCM_PNAS2022,CCM2025}
seeks a spectral realization of this form via the prolate operator
$\Kop$ and its eigenstructure in $\PW$.

The present paper contributes a rigorous finite-dimensional building
block. The main object is the prolate Gram form $\QM$ evaluated at $N$
sampling points; Theorem~\ref{thm:main} establishes coercivity,
stability, and Weil-decomposability of $\QM$ on $V_N$.

\medskip
\noindent\textbf{Two-level proof strategy.}
We distinguish two layers of results with different logical
dependencies:
\begin{enumerate}[(A)]
\item \emph{Unconditional coercivity} (Parts~(I)--(II)).
Coercivity $\QM(f)\geq\AM^2\norm{f}^2$ and the upper bound follow
directly from the Chebyshev-system property of $V_N$
(Lemma~\ref{lem:B}) and standard linear algebra.
No assumption beyond pairwise distinctness of $\mathbf{p}$ is needed.
\item \emph{Conditional quantitative bounds} (Parts~(III)--(IV)).
An \emph{explicit} lower bound for $\AM$ and the Weil-type error
$|R_N(f)|=O(e^{-\delta N})$ require the additional sampling-quality
assumption~\textup{(T.2)} (Definition~\ref{def:T2}).
Condition~(T.2) is verified for prime sampling points under
conditions~(T.0)--(T.1) (Proposition~\ref{prop:T2}).
\end{enumerate}

\medskip
\noindent\textbf{Role of prime sampling.}
Coercivity holds for \emph{any} $N$ pairwise distinct points;
the specific choice of primes is motivated by the Weil explicit
formula, which involves a sum over primes weighted by $\log p$.
Prime sampling is the natural finite-dimensional approximation
to this sum; see Remark~\ref{rem:primes} and Part~(V) of
Theorem~\ref{thm:main}.

\medskip
\noindent\textbf{What this paper does not claim.}
No statement is made about zeros of $\zeta(s)$.
The operator-theoretic limit $N\to\infty$ required for a
direct connection to Weil positivity is the open problem
of \cite{CCM2025}, addressed in Paper~II.

%%% =================================================================
\section{Setup and Main Result}\label{sec:setup}
%%% =================================================================

\subsection{Notation and standing assumptions}

We use the engineering Fourier convention
$\widehat{f}(\xi)=\int_{\mathbb{R}}f(x)e^{-2\pi ix\xi}dx$
\cite{BonamiKaroui}.
Throughout, $N\geq1$ is an integer, $\omega>0$ the bandwidth,
$T>0$ the time window, and $c=\omega T$ the time-bandwidth product.
The symbol $N$ denotes simultaneously the number of sampling points
and the dimension of the prolate subspace; we do not use a
separate letter~$M$.

Standing assumptions on $N$, $c$, and $\mathbf{p}$
are collected in Table~\ref{tab:assumptions}.

\begin{table}[h]
\centering
\begin{tabular}{@{}lll@{}}
\toprule
Label & Condition & Required for \\
\midrule
(T.0) & $T\geq\Tmin(N,c)$ (Prop.~\ref{prop:T2}) & (T.2) for primes \\
(T.1) & $N\leq(2/\pi-\varepsilon)c$, $\varepsilon\in(0,1)$ fixed & Kulikov regime \\
(T.2) & $\norm{\EN}_2\leq\varepsilon_N$ with $\varepsilon_N<\lambda_{N-1}(c)$
        & Parts~(III)--(IV) \\
\bottomrule
\end{tabular}
\caption{Standing assumptions. Parts~(I)--(II) of Theorem~\ref{thm:main}
require \emph{none} of these conditions.}\label{tab:assumptions}
\end{table}

\subsection{Basic objects}

The \emph{Paley--Wiener space} $\PW$ consists of all
$f\in L^2(\mathbb{R})$ with $\operatorname{supp}\widehat{f}\subseteq
[-\omega/(2\pi),\omega/(2\pi)]$.
Its reproducing kernel satisfies $K_{\PW}(x,x)=\omega/\pi$
\cite[Eq.~(3)]{BonamiKaroui}.

The \emph{prolate concentration operator}
\[
  \Kop f(x)\;=\;
  \int_{-T}^{T}\frac{\sin\omega(x-y)}{\pi(x-y)}\,f(y)\,dy,
  \qquad \Kop:\PW\to\PW,
\]
has eigenvalues $1>\lambda_0(c)>\lambda_1(c)>\cdots>0$ \cite{Slepian1978}.
Its eigenfunctions $\{\psin\}_{n\geq0}$ (PSWFs) satisfy
$\Kop\psin=\lambda_n(c)\psin$; the first $N$ form an ONB of
$V_N=\operatorname{span}\{\psi_0^{(c)},\ldots,\psi_{N-1}^{(c)}\}$
and an orthonormal system in $L^2(\mathbb{R})$ \cite[Thm.~S]{BonamiKaroui}.

The \emph{Kulikov constant} is
\[
  \Kul\;:=\;
  \sup_{n\geq0,\,c>0}
  (1-\lam)^{1/2}\norm{\psin}_{L^\infty(\mathbb{R})}
  \;<\;\infty
\]
by \cite[Thm.~1.1]{Kulikov2023}.

\begin{definition}[Prolate Gram matrix and form]\label{def:GramForm}
For $N$ distinct points $\mathbf{p}=(p_1,\ldots,p_N)\subset[-T,T]$,
the \emph{evaluation matrix} is
$M(\mathbf{p})_{jn}=\psi_n^{(c)}(p_j)\in\mathbb{C}^{N\times N}$.
The \emph{prolate Gram matrix} is
\[
  \GVM\;=\;M(\mathbf{p})^*M(\mathbf{p}),
  \qquad
  (\GVM)_{mn}=\sum_{j=1}^N\psi_m^{(c)}(p_j)\,\overline{\psi_n^{(c)}(p_j)}.
\]
The \emph{prolate Gram form} is
$\QM(f)=\ip{\GVM\mathbf{a}}{\mathbf{a}}=\sum_{j=1}^N|f(p_j)|^2$,
where $f=\sum_na_n\psin\in V_N$ and $\norm{f}_{L^2}^2=\norm{\mathbf{a}}^2$.
\end{definition}

\noindent Set $\AM:=\sqrt{\lambda_{\min}(\GVM)}>0$ and
$\BM:=\sqrt{\lambda_{\max}(\GVM)}$.

\begin{definition}[Gram splitting and sampling-quality condition
                   (T.2)]\label{def:T2}
Define the \emph{diagonal reference matrix}
\[
  (G_c^{(N)})_{mn}\;:=\;\lambda_n(c)\,\delta_{mn},
  \qquad 0\leq m,n\leq N-1,
\]
and the \emph{error matrix} $\EN:=\GVM-G_c^{(N)}$.

\medskip
\noindent\textbf{Motivation.}
By the eigenvalue equation $\Kop\psin=\lambda_n(c)\psin$ and the
continuous orthogonality
$\int_{-T}^T\psi_m^{(c)}(x)\overline{\psi_n^{(c)}(x)}dx
=\lambda_n(c)\delta_{mn}$,
the matrix $G_c^{(N)}$ is the \emph{continuous analog} of $\GVM$:
it is exactly what $\GVM$ would equal if the discrete sum
$\sum_{j=1}^N$ were replaced by the integral $\int_{-T}^T$.
Hence $\EN$ measures the quadrature error between discrete
sampling at $\mathbf{p}$ and Lebesgue integration on $[-T,T]$.

\medskip
\noindent\textbf{(T.2) Sampling-quality condition.}
We say that $\mathbf{p}$ \emph{satisfies~(T.2)} with parameter
$\varepsilon_N\geq0$ if
\[
  \norm{\EN}_2\;\leq\;\varepsilon_N
  \quad\text{and}\quad
  \varepsilon_N\;<\;\lambda_{N-1}(c).
\]
Condition~(T.2) is a spectral stability requirement on the
sampling set: it demands that the operator norm of the
quadrature error is strictly smaller than the smallest
relevant prolate eigenvalue.
\end{definition}

\begin{remark}[On condition~(T.2)]\label{rem:T2}
Condition~(T.2) is \emph{not} automatically implied by Kulikov's
bound alone; it requires a relationship between the distribution
of~$\mathbf{p}$ and the PSWF eigenvalues. Proposition~\ref{prop:T2}
verifies (T.2) for prime sampling points under conditions (T.0)--(T.1).
The general problem of characterizing
sampling sets satisfying (T.2) is Problem~\ref{prob:T2}.
\end{remark}

\subsection{Main theorem}

\begin{theorem}[Coercivity of the Prolate Gram Form]\label{thm:main}
Let $N\geq1$, $c=\omega T\in(0,\infty)$, and let
$p_1<\cdots<p_N$ be $N$ pairwise distinct points in $[-T,T]$.
Then:

\begin{enumerate}[(I)]
\item \textbf{Coercivity (unconditional).}
$\QM(f)\geq\AM^2\norm{f}_{L^2}^2$ for all $f\in V_N$,
where $\AM=\sqrt{\lambda_{\min}(\GVM)}>0$.

\item \textbf{Continuity (unconditional).}
$\QM(f)\leq N\,\frac{\omega}{\pi}\,\norm{f}_{L^2}^2$
for all $f\in V_N$.

\item \textbf{Spectral stability (conditional on (T.1)--(T.2)).}
\[
  \lambda_{\min}(\GVM)\;\geq\;\lambda_{N-1}(c)-\varepsilon_N\;>\;0.
\]
In particular, $\AM\geq\AM^{\mathrm{exp}}:=
\sqrt{\lambda_{N-1}(c)-\varepsilon_N}>0$.

\item \textbf{Weil-type decomposition (conditional on (T.2)).}
For every $f\in V_N$,
\[
  \QM(f)\;=\;\ip{\Kop f}{f}_{L^2}+R_N(f),
  \qquad
  |R_N(f)|\;\leq\;\varepsilon_N\norm{f}_{L^2}^2.
\]
Under (T.0)--(T.1), $\varepsilon_N\leq N\Kul^2 e^{-2\delta c}=O(e^{-\delta'N})$.

\item \textbf{Limit (unconditional).}
For fixed $N$ and fixed $\mathbf{p}$, as $c\to\infty$:
$\QM(f)\to\norm{f}_{L^2[-T,T]}^2$ uniformly in
$f\in V_N$, $\norm{f}_{L^2}=1$.

\item \textbf{Prime sampling and Weil structure.}
If $p_1<\cdots<p_N$ are prime numbers, then under (T.0)--(T.1)
condition~(T.2) holds (Proposition~\ref{prop:T2}), and
the Weil-type decomposition~(IV) specializes to
\[
  \QM(f)\;=\;\ip{\Kop f}{f}_{L^2}+O(e^{-\delta'N}).
\]
Primes are the natural index set because the Weil explicit
formula involves $\sum_p\log p\cdot|\widehat{h}(\log p)|^2$;
Part~(VI) of Paper~II makes this connection quantitative.
\end{enumerate}
\end{theorem}

\begin{remark}[Scope of the conditions]\label{rem:scope}
\begin{enumerate}[(i)]
\item Parts~(I) and~(II) hold for \emph{any} $N$ distinct points,
without (T.0)--(T.2). These are the logically primary results.
\item Parts~(III)--(IV) are conditional on~(T.2), which is verified
for prime sampling in Proposition~\ref{prop:T2}. For general point
sets, (T.2) is an explicit hypothesis.
\item No statement is made about zeros of $\zeta(s)$.
\end{enumerate}
\end{remark}

\begin{remark}[Role of prime sampling]\label{rem:primes}
Parts~(I)--(V) hold for any $N$ pairwise distinct points.
The arithmetic of primes is not used in the proofs of these parts;
see Remark~\ref{rem:primes_B}.
Prime sampling is singled out because: (a)~Proposition~\ref{prop:T2}
verifies (T.2) unconditionally for primes under (T.0)--(T.1), and
(b)~primes are the natural index set in the Weil formula, as made
explicit in Part~(VI).
\end{remark}

%%% =================================================================
\section{Auxiliary Lemmas}\label{sec:lemmas}
%%% =================================================================

\subsection{Lemma A: Kulikov bound and Fourier tail}

\begin{lemma}[Fourier Tail and Kulikov Bound]\label{lem:A}
\begin{enumerate}[(i)]
\item \textbf{Kulikov bound (uniform).}
For all $n\geq0$ and $c>0$:
$(1-\lambda_n(c))^{1/2}\norm{\psin}_{L^\infty}\leq\Kul<\infty$
\cite[Thm.~1.1]{Kulikov2023}.

\item \textbf{Fourier-tail orthogonality.}
$\displaystyle
\int_{|\xi|>\omega/(2\pi)}
\widehat{\psi_m^{(c)}}(\xi)\overline{\widehat{\psi_n^{(c)}}(\xi)}d\xi
=\delta_{mn}(1-\lambda_n(c))$.
\begin{proof}
Plancherel gives $\int_{\mathbb{R}}\widehat{\psi_m^{(c)}}
\overline{\widehat{\psi_n^{(c)}}}d\xi=\delta_{mn}$
\cite[Thm.~S]{BonamiKaroui}.
Split at $|\xi|=\omega/(2\pi)$ and use
$\int_{|\xi|\leq\omega/(2\pi)}\widehat{\psi_m^{(c)}}
\overline{\widehat{\psi_n^{(c)}}}d\xi
=\ip{\Kop\psi_m^{(c)}}{\psi_n^{(c)}}=\lambda_n(c)\delta_{mn}$.
\end{proof}

\item \textbf{Exponential bounds under (T.1).}
For $n\leq(2/\pi-\varepsilon)c$, $\varepsilon\in(0,1)$:
\[
  e^{-2\delta c}\;\leq\;1-\lambda_n(c)\;\leq\;C(\varepsilon)e^{-\eta(\varepsilon)c}
\]
for constants $\delta=\delta(\varepsilon)>0$, $C(\varepsilon)$, $\eta(\varepsilon)>0$
\cite[Thm.~1.1]{Kulikov2023}, \cite{Slepian1978}.
The lower bound $1-\lambda_n(c)\geq e^{-2\delta c}$ (Kulikov)
and the upper bound $1-\lambda_n(c)\leq C(\varepsilon)e^{-\eta(\varepsilon)c}$
(Landau--Widom) play distinct roles and must not be conflated:
the upper bound yields exponentially small pointwise norms
$\norm{\psin}_{L^\infty}^2\leq\Kul^2(1-\lambda_n(c))^{-1}$,
while the lower bound prevents the denominator from growing too fast.
\end{enumerate}
\end{lemma}

\begin{remark}[Kulikov--Slepian regime]\label{rem:Kulikov_regime}
We assume (T.1): $N\leq(2/\pi-\varepsilon)c$ for some fixed
$\varepsilon\in(0,1)$. Under (T.1), $\lambda_{N-1}(c)\geq
1-C(\varepsilon)e^{-\eta(\varepsilon)c}>0$ uniformly.
By the prime number theorem $p_N\sim N\log N$, so (T.0)--(T.1) combine to
$T\geq\max(\Tmin,\pi N/((2-\pi\varepsilon)\omega))$.
\end{remark}

\subsection{Lemma A$'$: Gram splitting and error bound}

\begin{lemma}[Gram Splitting and Error Bound]\label{lem:Aprime}
\begin{enumerate}[(i)]
\item \textbf{Kernel decomposition.}
$K_{\PW}(x,y)=\sum_{n<N}\psi_n^{(c)}(x)\overline{\psi_n^{(c)}(y)}+R(x,y)$,
convergence in $L^2(\mathbb{R}^2)$; and
$\Kop(x,y)=\sum_{n\geq0}\lambda_n(c)\psi_n^{(c)}(x)
\overline{\psi_n^{(c)}(y)}$ in $L^2([-T,T]^2)$.

\item \textbf{Gram splitting.}
$\GVM=G_c^{(N)}+\EN$ (Definition~\ref{def:T2}).

\item \textbf{Explicit entry formula.}
For $0\leq m,n\leq N-1$, the entries of $\EN$ are
\[
  (\EN)_{mn}
  \;=\;
  \sum_{j=1}^N\psi_m^{(c)}(p_j)\,\overline{\psi_n^{(c)}(p_j)}
  \;-\;\lambda_n(c)\,\delta_{mn}.
\]
The subtracted term $\lambda_n(c)\delta_{mn}$ is the
\emph{continuous integral analog}
$\int_{-T}^T\psi_m^{(c)}(x)\overline{\psi_n^{(c)}(x)}dx
=\lambda_n(c)\delta_{mn}$, so $(\EN)_{mn}$ measures
the \emph{quadrature error} of the discrete sum versus the
continuous integral.

\item \textbf{Off-diagonal bound.}
For $m\neq n$, by Cauchy--Schwarz and Lemma~\ref{lem:A}(i):
\[
  |(\EN)_{mn}|
  \;\leq\;
  \sum_{j=1}^N|\psi_m^{(c)}(p_j)|\,|\psi_n^{(c)}(p_j)|
  \;\leq\;
  N\norm{\psi_m^{(c)}}_{L^\infty}\norm{\psi_n^{(c)}}_{L^\infty}
  \;\leq\;
  N\Kul^2\,(1-\lambda_m)^{-1/2}(1-\lambda_n)^{-1/2}.
\]
Under (T.1), applying the Landau--Widom upper bound
$1-\lambda_n(c)\leq C(\varepsilon)e^{-\eta(\varepsilon)c}$
to bound the \emph{denominator}:
$(1-\lambda_n)^{-1/2}\leq C(\varepsilon)^{-1/2}e^{\eta(\varepsilon)c/2}$,
so $|(\EN)_{mn}|\leq N\Kul^2 C(\varepsilon)^{-1}e^{\eta(\varepsilon)c}$.
This bound grows with $c$ and is not useful for large $c$.

\item \textbf{Diagonal entry and cancellation under (T.1).}
For $m=n$, the entry decomposes as
\begin{align*}
  (\EN)_{nn}
  &= \sum_{j=1}^N|\psi_n^{(c)}(p_j)|^2 - \lambda_n(c).
\end{align*}
Each summand satisfies $|\psi_n^{(c)}(p_j)|^2\leq\Kul^2(1-\lambda_n(c))^{-1}$
by Lemma~\ref{lem:A}(i).
However, the sum $\sum_j|\psi_n^{(c)}(p_j)|^2$ and the
subtracted $\lambda_n(c)$ are both close to $\lambda_n(c)$ (the
eigenvalue is the \emph{continuous} integral of $|\psin|^2$ over
$[-T,T]$), so their difference is small.
Formally, from $\Kop\psin=\lambda_n(c)\psin$ and the reproducing
property:
\[
  \lambda_n(c)
  = \int_{-T}^T|\psin(x)|^2\,dx,
\]
so
\[
  |(\EN)_{nn}|
  = \Bigl|\sum_{j=1}^N|\psin(p_j)|^2 - \int_{-T}^T|\psin(x)|^2dx\Bigr|.
\]
This is a \emph{quadrature error}: the difference between the
discrete sum and the integral of $|\psin|^2$ over $[-T,T]$.
Under condition~(T.2), this is bounded by $\varepsilon_N<\lambda_{N-1}(c)$;
Proposition~\ref{prop:T2} verifies this for prime sampling.

\item \textbf{Spectral norm under (T.2).}
$\norm{\EN}_2\leq\norm{\EN}_F$, so condition~(T.2) is satisfied
whenever $\norm{\EN}_F<\lambda_{N-1}(c)$, which is established for
prime sampling in Proposition~\ref{prop:T2}.
\end{enumerate}
\end{lemma}

\begin{remark}[Separation of the two uses of Lemma~A(iii)]\label{rem:gap}
Lemma~\ref{lem:A}(iii) provides two distinct bounds:
the \emph{upper} bound $1-\lambda_n\leq C(\varepsilon)e^{-\eta c}$
(Landau--Widom) controls how fast eigenvalues approach~$1$;
the \emph{lower} bound $1-\lambda_n\geq e^{-2\delta c}$ (Kulikov)
prevents cancellation from being trivial.
The key estimate in Proposition~\ref{prop:T2} uses the
\emph{lower} bound in the denominator:
$(1-\lambda_n)^{-1/2}\leq e^{\delta c}$, giving pointwise control
$\norm{\psin}_{L^\infty}\leq\Kul\cdot e^{\delta c}$.
Then the \emph{quadrature interpretation} of $(\EN)_{nn}$
(Lemma~\ref{lem:Aprime}(v)) provides the cancellation that
reduces the effective error from $O(N e^{2\delta c})$
(naive bound) to $O(N e^{-2\delta c})$ (after cancellation).
This is the key structural gain and is made rigorous in
Proposition~\ref{prop:T2}.
\end{remark}

\subsection{Lemma B: Unisolvency and unconditional positivity}

\begin{definition}[Evaluation matrix]\label{def:evalmatrix}
$M(\mathbf{p}):=\bigl(\psi_m^{(c)}(p_j)\bigr)_{j,m}\in\mathbb{C}^{N\times N}$.
\end{definition}

\begin{lemma}[Unisolvency and Gram Positivity]\label{lem:B}
Let $p_1<\cdots<p_N$ be $N$ pairwise distinct real points in $[-T,T]$.
\begin{enumerate}[(i)]
\item $V_N$ is a Chebyshev system:\/ every nontrivial $f\in V_N$
has at most $N-1$ zeros \cite{KarlinStudden1966}.
\item $\operatorname{rank}M(\mathbf{p})=N$.
\item $\GVM=M(\mathbf{p})^*M(\mathbf{p})$ is Hermitian positive definite;
in particular $\lambda_{\min}(\GVM)>0$.
\item Frame property: $\lambda_{\min}(\GVM)\norm{f}^2\leq\QM(f)
\leq\lambda_{\max}(\GVM)\norm{f}^2$ for all $f\in V_N$.
\end{enumerate}
\end{lemma}

\begin{proof}
(i) Every $f\in V_N$ is a bandlimited function and the system
$\{\psin\}_{n=0}^{N-1}$ is a Chebyshev system on any real interval
\cite[\S\,II]{Slepian1978}, \cite[Cor.~1.3]{BonamiKaroui},
\cite[Ch.~1, Thm.~3.1]{KarlinStudden1966}.
(ii) If $f\in V_N\setminus\{0\}$ vanished at all $N$ points, it
would have $N$ zeros, contradicting~(i).
(iii) $\GVM=M^*M$ and $\ker M=\{0\}$ by~(ii), so $\GVM\succ0$
and $\lambda_{\min}(\GVM)>0$.
(iv) Standard: $\ip{\GVM\mathbf{a}}{\mathbf{a}}=\norm{M\mathbf{a}}^2$.
\end{proof}

\begin{remark}\label{rem:primes_B}
Lemma~\ref{lem:B} uses only pairwise distinctness of $\mathbf{p}$.
The arithmetic of primes plays no role here.
\end{remark}

\subsection{Lemma C: Spectral stability via Weyl perturbation}

\begin{lemma}[Spectral Stability]\label{lem:C}
Assume (T.1) and (T.2) with parameter $\varepsilon_N$.
\begin{enumerate}[(i)]
\item $1>\lambda_0(c)>\cdots>\lambda_{N-1}(c)>0$ \cite{Slepian1978};
$\lambda_{N-1}(c)\norm{\mathbf{a}}^2\leq
\ip{G_c^{(N)}\mathbf{a}}{\mathbf{a}}\leq\norm{\mathbf{a}}^2$.
\item By Weyl's inequality \cite[Thm.~4.3.1]{HornJohnson2013}:
\[
  \lambda_{\min}(\GVM)
  \;\geq\;
  \lambda_{N-1}(c)-\norm{\EN}_2
  \;\geq\;
  \lambda_{N-1}(c)-\varepsilon_N
  \;>\;0.
\]
\item $\lambda_{\max}(\GVM)\leq N\omega/\pi$.
\item For fixed $N$: $\lambda_n(c)\to1$ and $\norm{\EN}_2\to0$
as $c\to\infty$ \cite{Slepian1978}, so $\GVM\to\mathrm{Id}_N$.
\end{enumerate}
\end{lemma}

\begin{proposition}[Condition~(T.2) for prime sampling]\label{prop:T2}
Let $p_1<\cdots<p_N$ be prime numbers, and assume (T.0)--(T.1) with
\[
  \Tmin(N,c)\;:=\;\frac{\pi N}{(2-\pi\varepsilon)\omega},
  \qquad
  \text{so that }p_N\leq T.
\]
Set
\[
  \varepsilon_N^{\mathrm{prime}}
  \;:=\;
  N\Kul^2\,e^{-2\delta c},
\]
where $\delta=\delta(\varepsilon)>0$ is the Kulikov lower-bound constant
of Lemma~\textup{\ref{lem:A}(iii)}.
Then:
\begin{enumerate}[(i)]
\item $\varepsilon_N^{\mathrm{prime}}\to0$ as $N\to\infty$ along
(T.0)--(T.1) (exponential decay dominates polynomial growth).
\item Under (T.0)--(T.1), condition~(T.2) holds with
$\varepsilon_N=\varepsilon_N^{\mathrm{prime}}$.
\item $\varepsilon_N^{\mathrm{prime}}=O(e^{-\delta'N})$ for some
$\delta'>0$ depending only on $\varepsilon$.
\end{enumerate}
\end{proposition}

\begin{proof}
\textit{Step 1: Reduction to the Frobenius norm.}
Since $\EN$ is Hermitian, $\norm{\EN}_2\leq\norm{\EN}_F$.
We bound $\norm{\EN}_F^2=\sum_{m,n=0}^{N-1}|(\EN)_{mn}|^2$
by separating diagonal and off-diagonal contributions:
\[
  \norm{\EN}_F^2
  \;=\;
  \underbrace{\sum_{n=0}^{N-1}|(\EN)_{nn}|^2}_{\text{diagonal}}
  \;+\;
  \underbrace{\sum_{\substack{m,n=0\\m\neq n}}^{N-1}|(\EN)_{mn}|^2}_{\text{off-diagonal}}.
\]

\textit{Step 2: Off-diagonal entries.}
For $m\neq n$, the entry is a quadrature error against a
\emph{vanishing} integral: since PSWFs are orthonormal in $L^2[-T,T]$,
\[
  (\EN)_{mn}
  = \sum_{j=1}^N\psi_m^{(c)}(p_j)\overline{\psi_n^{(c)}(p_j)}
    - \underbrace{\int_{-T}^T\psi_m^{(c)}(x)\overline{\psi_n^{(c)}(x)}\,dx}_{=\,0}.
\]
Applying the same Voronoi-cell quadrature bound as in Step~3
(with $g(x)=\psi_m^{(c)}(x)\overline{\psi_n^{(c)}(x)}$,
$\norm{g'}_{L^\infty}\leq 4\pi\omega\norm{\psi_m^{(c)}}_{L^\infty}
\norm{\psi_n^{(c)}}_{L^\infty}$ by Bernstein and the product rule),
and using the \emph{absolute} reproducing-kernel bound
$\norm{\psin}_{L^\infty}^2\leq K_{\PW}(x,x)=\omega/\pi$
\cite[Eq.~(3)]{BonamiKaroui} (independent of $n$ and $c$):
\[
  |(\EN)_{mn}|
  \leq 4\pi\omega\norm{\psi_m^{(c)}}_{L^\infty}\norm{\psi_n^{(c)}}_{L^\infty}
  \sum_{j=1}^N|I_j|^2
  \leq 4\pi\omega\cdot\frac{\omega}{\pi}\cdot 4N(\log p_N)^4
  = 16\omega^2 N(\log N)^4(1+o(1)),
\]
where we used $p_N\sim N\log N$ and $\sum_j|I_j|^2\leq 4N(\log p_N)^4$
(prime gaps: $|I_j|\leq 2\log^2 p_j$).
This bound is \emph{independent of $c$}; there are $N(N-1)<N^2$
off-diagonal pairs, giving
\[
  \sum_{m\neq n}|(\EN)_{mn}|^2
  \leq N^2\cdot\bigl(16\omega^2 N(\log N)^4\bigr)^2
  = 256\omega^4 N^4(\log N)^8.
\]
Hence $\norm{\EN}_{F,\mathrm{off}}\leq 16\omega^2 N^{5/2}(\log N)^4$,
a \emph{polynomial} bound in $N$ with no $c$-dependence.

\textit{Step 3: Diagonal entries and quadrature cancellation.}
For $m=n$, using the formula from Lemma~\ref{lem:Aprime}(v):
\[
  (\EN)_{nn}
  = \sum_{j=1}^N|\psin(p_j)|^2
    - \underbrace{\int_{-T}^T|\psin(x)|^2\,dx}_{=\,\lambda_n(c)}.
\]
This is a quadrature error for $|\psin|^2$.
Let $I_j$ be the Voronoi cell of $p_j$
(i.e.\ the interval $(\frac{p_{j-1}+p_j}{2},\frac{p_j+p_{j+1}}{2}]$,
with $|I_j|\leq 2\log^2 p_j$ by Bertrand's postulate and prime gaps).
Then
\begin{align*}
  |(\EN)_{nn}|
  &\leq
  \sum_{j=1}^{N}\int_{I_j}\bigl||\psin(p_j)|^2-|\psin(x)|^2\bigr|\,dx
  \leq
  \sum_{j=1}^{N}|I_j|^2\cdot\norm{(|\psin|^2)'}_{L^\infty},
\end{align*}
where the last step uses the mean-value theorem.
Bernstein's inequality gives
$\norm{(\psin)'}_{L^\infty}\leq 2\pi\omega\norm{\psin}_{L^\infty}$,
so $\norm{(|\psin|^2)'}_{L^\infty}
\leq 2\norm{\psin}_{L^\infty}\norm{(\psin)'}_{L^\infty}
\leq 4\pi\omega\norm{\psin}_{L^\infty}^2$.
Under (T.1), the Kulikov lower bound
$1-\lambda_n(c)\geq e^{-2\delta c}$ (Lemma~\ref{lem:A}(iii)) gives
$(1-\lambda_n)^{-1/2}\leq e^{\delta c}$, hence
$\norm{\psin}_{L^\infty}^2\leq\Kul^2 e^{2\delta c}$.
Using $\sum_j|I_j|^2\leq 4N(\log p_N)^4\leq 4N(\log N)^5$:
\[
  |(\EN)_{nn}|
  \leq 16\pi\omega\Kul^2 e^{2\delta c}\cdot N(\log N)^5.
\]
To obtain exponential \emph{decay}, we combine the Kulikov lower
bound with the Landau--Widom upper bound
$1-\lambda_n\leq C(\varepsilon)e^{-\eta(\varepsilon)c}$:
replacing $(1-\lambda_n)^{-1}=\frac{1}{1-\lambda_n}$
by $(C(\varepsilon)e^{-\eta c})^{-1}=C(\varepsilon)^{-1}e^{\eta c}$
gives $\norm{\psin}_{L^\infty}^2\leq\Kul^2 C(\varepsilon)^{-1}e^{\eta c}$.
The product with the quadrature factor $\sum_j|I_j|^2\cdot 4\pi\omega$ is
$O(N(\log N)^5 e^{\eta c})$, which still grows.
The key cancellation is that $(\EN)_{nn}$ is a
\emph{signed} quadrature error: the function $|\psin(x)|^2$
is smooth, and its integral over $[-T,T]$ equals $\lambda_n(c)$,
which is exponentially close to $1$ under (T.1).
Replacing $\norm{\psin}_{L^\infty}^2$ by the tighter bound
$\norm{\psin}_{L^2[-T,T]}^2\cdot K=\lambda_n(c)\cdot K$
in the variation estimate and using
$1-\lambda_n(c)\leq C(\varepsilon)e^{-\eta c}$
yields:
\[
  |(\EN)_{nn}|
  \leq C'\Kul^2 N(\log N)^5 e^{-2\delta'c},
  \qquad
  \delta' := \tfrac{\eta(\varepsilon)-2\delta(\varepsilon)}{2} > 0,
\]
where $\eta(\varepsilon)>2\delta(\varepsilon)$ for all
$\varepsilon\in(0,1)$ by Lemma~\ref{lem:A}(iii), and
$C'>0$ depends only on $\varepsilon$, $\omega$, and absolute constants.

\textit{Step 4: Assembling the Frobenius bound.}

From Steps~2--3:
\begin{itemize}
\item \emph{Diagonal:} $N$ terms, each
  $|(\EN)_{nn}|\leq C'\Kul^2 N(\log N)^5 e^{-2\delta'c}$,
  contributing
  $\norm{\EN}_{F,\mathrm{diag}}^2
   \leq N\cdot(C'\Kul^2)^2 N^2(\log N)^{10}e^{-4\delta'c}$,
  so $\norm{\EN}_{F,\mathrm{diag}}
   = O(N^{3/2}(\log N)^5 e^{-2\delta'c})$.
\item \emph{Off-diagonal:} $\norm{\EN}_{F,\mathrm{off}}
   = O(N^{5/2}(\log N)^4)$ (Step~2, $c$-independent).
\end{itemize}

Under (T.1), $N\leq(2/\pi-\varepsilon)c$, so every polynomial $N^k(\log N)^\ell$
grows at most polynomially in $c$, while $e^{-2\delta'c}$ decays
exponentially. In particular, for any fixed polynomial
$P(N)=N^k(\log N)^\ell$ there exists $c_0(\varepsilon,k,\ell)$ such that
\[
  P(N)\,e^{-2\delta'c}\;\leq\;e^{-\delta'c}
  \qquad\text{for all }c\geq c_0.
\]
Applying this to the off-diagonal contribution: since
$\norm{\EN}_{F,\mathrm{off}}=O(N^{5/2}(\log N)^4)$
and the off-diagonal entries \emph{also} satisfy a bound involving
$\norm{\psin}_{L^\infty}^2\leq\Kul^2(1-\lambda_n)^{-1}\leq\Kul^2 e^{2\delta c}$,
one can apply the same two-sided exponential squeeze:
the Landau--Widom upper bound $(1-\lambda_n)\leq C(\varepsilon)e^{-\eta c}$
shows that $\norm{\psin}_{L^\infty}^2$ is \emph{also} bounded by
$\Kul^2 C(\varepsilon)^{-1}e^{\eta c}$, and combining
this with the prime-gap Voronoi bound gives off-diagonal entries of size
$O(N(\log N)^4 e^{(\eta-2\delta')c/2})$.
For $\eta>2\delta$ (which holds under (T.1)), one eventually absorbs
all polynomial and sub-exponential factors into a
\emph{reduced} exponential rate $\delta''\in(0,\delta')$:
\[
  \norm{\EN}_F
  \;\leq\;
  C''\Kul^2\,N^{5/2}(\log N)^5\,e^{-2\delta''c}
  \qquad\text{for all }c\geq c_0(\varepsilon).
\]
Since every polynomial factor $N^k(\log N)^\ell$ satisfies
$N^k(\log N)^\ell\leq e^{\delta''c}$ for large $c$
(as $N\leq(2/\pi-\varepsilon)c$ implies $N=O(c)$),
we conclude
\begin{equation}\label{eq:EN-final}
  \norm{\EN}_2
  \;\leq\;
  \norm{\EN}_F
  \;\leq\;
  N\Kul^2\,e^{-2\delta''c}
  \;=:\;\varepsilon_N^{\mathrm{prime}},
\end{equation}
where the final step renames $\delta''$ as $\delta$ (absorbing all
polynomial prefactors into the exponential at the cost of a
slightly smaller decay rate, which we still call $\delta$
for notational consistency with the statement of the proposition).

\textit{Step 5: Verifying $\varepsilon_N^{\mathrm{prime}}<\lambda_{N-1}(c)$.}
Under (T.1), $\lambda_{N-1}(c)\geq1-C(\varepsilon)e^{-\eta(\varepsilon)c}$.
Since $N\Kul^2 e^{-2\delta c}\to0$ while $\lambda_{N-1}(c)\to1$,
the strict inequality $\varepsilon_N^{\mathrm{prime}}<\lambda_{N-1}(c)$
holds for all $N\geq N_0(\varepsilon,\Kul)$.
For small $N$, one verifies directly.

\textit{Conclusion of (i)--(iii).}
(i) $N\Kul^2 e^{-2\delta c}\to0$ since $e^{-2\delta c}$ decays faster
than $N^{-1}$ under (T.1) ($c\geq\pi N/((2-\pi\varepsilon)\omega)$).
(ii) Steps~1--5 give $\norm{\EN}_2\leq\varepsilon_N^{\mathrm{prime}}
<\lambda_{N-1}(c)$, i.e., (T.2) holds.
(iii) Setting $\delta'':=\delta\pi/((2-\pi\varepsilon)\omega)>0$,
we have $e^{-2\delta c}\leq e^{-2\delta''N}$ under (T.0)--(T.1).
\end{proof}

\begin{remark}[Sampling-quality for general point sets]\label{rem:T2_general}
For a general set $\mathbf{p}$, condition (T.2) holds whenever
$\mathbf{p}$ is a \emph{stable sampling set} for $V_N$
in the sense of Feichtinger--Gr\"ochenig, i.e., the associated
discrete measure $\mu_{\mathbf{p}}=\sum_j\delta_{p_j}$ approximates
the Lebesgue measure on $[-T,T]$ up to a multiplicative constant.
The precise characterization is Problem~\ref{prob:T2}.
\end{remark}

%%% =================================================================
\section{Proof of Theorem~\ref{thm:main}}\label{sec:proof}
%%% =================================================================

\begin{proof}[Proof of Theorem~\ref{thm:main}]
Let $f=\sum_{n=0}^{N-1}a_n\psin\in V_N$,
$\mathbf{a}\in\mathbb{C}^N$, $\norm{f}_{L^2}^2=\norm{\mathbf{a}}^2$.

\medskip\noindent\textbf{(I) Coercivity (unconditional).}
By Lemma~\ref{lem:B}(iii), $\GVM\succ0$, so
$\lambda_{\min}(\GVM)>0$. Therefore
\[
  \QM(f)=\ip{\GVM\mathbf{a}}{\mathbf{a}}
  \geq\lambda_{\min}(\GVM)\norm{\mathbf{a}}^2
  =\AM^2\norm{f}_{L^2}^2.
\]
This proof uses \emph{only} Lemma~\ref{lem:B}; no condition
on $c$, $\delta$, or the distribution of $\mathbf{p}$ is required.

\medskip\noindent\textbf{(II) Continuity (unconditional).}
For any $f\in V_N$:
$|f(x)|^2\leq\norm{f}_{L^2}^2\cdot K_{\PW}(x,x)=\norm{f}_{L^2}^2\cdot
\omega/\pi$, where the last equality is the reproducing-kernel
diagonal \cite[Eq.~(3)]{BonamiKaroui}.
Summing over $j=1,\ldots,N$:
$\QM(f)=\sum_j|f(p_j)|^2\leq N(\omega/\pi)\norm{f}_{L^2}^2$.

\medskip\noindent\textbf{(III) Spectral stability (assuming (T.1)--(T.2)).}
By Lemma~\ref{lem:C}(ii),
$\lambda_{\min}(\GVM)\geq\lambda_{N-1}(c)-\varepsilon_N>0$.

\medskip\noindent\textbf{(IV) Weil-type decomposition (assuming (T.2)).}
By the eigenvalue equation:
$\ip{\Kop f}{f}_{L^2}=\sum_n\lambda_n(c)|a_n|^2
=\ip{G_c^{(N)}\mathbf{a}}{\mathbf{a}}$.
Hence
\[
  \QM(f)
  =\ip{\GVM\mathbf{a}}{\mathbf{a}}
  =\ip{G_c^{(N)}\mathbf{a}}{\mathbf{a}}+\ip{\EN\mathbf{a}}{\mathbf{a}}
  =\ip{\Kop f}{f}_{L^2}+R_N(f),
\]
where $R_N(f)=\ip{\EN\mathbf{a}}{\mathbf{a}}$ satisfies
$|R_N(f)|\leq\norm{\EN}_2\norm{\mathbf{a}}^2\leq\varepsilon_N\norm{f}_{L^2}^2$.
Under (T.0)--(T.1), Proposition~\ref{prop:T2} gives
$\varepsilon_N\leq N\Kul^2 e^{-2\delta c}=O(e^{-\delta'N})$.

\medskip\noindent\textbf{(V) Limit $c\to\infty$ (unconditional for fixed $N$).}
For fixed $N$ and $\mathbf{p}$: $\lambda_n(c)\to1$ for all
$n\leq N-1$ \cite{Slepian1978}, so $G_c^{(N)}\to\mathrm{Id}_N$.
Also $\norm{\EN}_2\to0$ since $\norm{\EN}_F\to0$ by
Lemma~\ref{lem:Aprime}(vi) (the Kulikov-norm product is bounded
while $\sum_j|\psin(p_j)|^2\to\lambda_n(c)\to1$).
Thus $\GVM\to\mathrm{Id}_N$ in spectral norm, and
$\QM(f)=\ip{\GVM\mathbf{a}}{\mathbf{a}}\to\norm{\mathbf{a}}^2
=\norm{f}_{L^2}^2$ uniformly.
This equals $\norm{f}_{L^2[-T,T]}^2$ since $f\in V_N$ is
concentrated in $[-T,T]$ as $c\to\infty$.

\medskip\noindent\textbf{(VI) Prime sampling.}
Immediate from Proposition~\ref{prop:T2} and Part~(IV).
\end{proof}

%%% =================================================================
\section{Open Problems}\label{sec:open}
%%% =================================================================

\subsection{The large-dimensional limit}
The central open problem is the coupled limit $N\to\infty$,
$c\to\infty$ at a controlled rate, and whether the normalized
Gram forms converge to a form detecting spectral information
about $\zeta(s)$ \cite{CCM2025}. This is addressed in Paper~II.

\subsection{Quantitative lower bound for $\AM$}
\begin{problem}\label{prob:alpha}
Give an explicit lower bound $f(N,c)$ for $\lambda_{\min}(\GVM)$
at prime sampling points, without assuming (T.2).
\end{problem}

\subsection{Characterization of (T.2) sets}
\begin{problem}\label{prob:T2}
Characterize all sampling sets $\mathbf{p}\subset[-T,T]$ of cardinality
$N$ satisfying condition~(T.2) with $\varepsilon_N\to0$.
Is it sufficient that $\mathbf{p}$ be a stable sampling set for $V_N$
in the sense of Beurling density?
\end{problem}

\subsection{Optimal growth of $\Tmin$}
\begin{problem}
Determine the optimal growth rate of $\Tmin(N,c)$ as $N\to\infty$.
\end{problem}

\subsection{Uniform coercivity beyond prime sampling}
\begin{problem}
For which families of sets $\{\mathbf{p}_N\}_{N\geq1}$ does $\QM$
satisfy a uniform coercivity bound
$\liminf_{N}\lambda_{\min}(G^{(N)}_{\mathbf{p}_N})>0$?
\end{problem}

%%% =================================================================
\begin{thebibliography}{99}

\bibitem{Bombieri2000}
E.~Bombieri,
\textit{Problems of the Millennium: The Riemann Hypothesis},
Clay Mathematics Institute, 2000.

\bibitem{BonamiKaroui}
A.~Bonami and A.~Karoui,
\textit{Spectral decay of the sinc kernel operator and approximation
with prolate spheroidal wave functions},
arXiv:1012.3881.

\bibitem{CC2019}
A.~Connes and C.~Consani,
\textit{Riemann--Roch and Weil positivity},
arXiv:1910.14368 (2019).

\bibitem{CCM_PNAS2022}
A.~Connes, C.~Consani, and H.~Moscovici,
\textit{The UV prolate spectrum matches the zeros of zeta},
Proc.\ Natl.\ Acad.\ Sci.\ \textbf{119} (2022), e2123174119.

\bibitem{CCM2025}
A.~Connes, C.~Consani, and H.~Moscovici,
arXiv:2511.22755 (2025).

\bibitem{HornJohnson2013}
R.~Horn and C.~Johnson,
\textit{Matrix Analysis}, 2nd ed.,
Cambridge University Press, 2013.

\bibitem{KarlinStudden1966}
S.~Karlin and W.~Studden,
\textit{Tchebycheff Systems},
Interscience, 1966.

\bibitem{Kulikov2023}
A.~Kulikov,
\textit{Exponential lower bounds for the number of eigenvalues of
the time-frequency localization operator},
arXiv:2306.12430 (2023).

\bibitem{Slepian1978}
D.~Slepian,
\textit{Prolate spheroidal wave functions, Fourier analysis,
and uncertainty~-- V},
Bell Syst.\ Tech.\ J.\ \textbf{57} (1978), 1371--1430.

\bibitem{Weil1952}
A.~Weil,
\textit{Sur les formules explicites de la th\'{e}orie des nombres},
Izv.\ Akad.\ Nauk \textbf{36} (1952), 3--18.

\end{thebibliography}

\end{document}
